\section{Simple Oracle Reductions}

We start by introducing a number of simple oracle reductions that serve as fundamental building blocks for more complex proof systems. These components can be composed together to construct larger protocols.

\subsection{Trivial Reduction}

The simplest possible oracle reduction is one that performs no computation at all. Both the prover and verifier simply pass their inputs through unchanged. While seemingly trivial, this reduction serves as an important identity element for composition and provides a base case for lifting and lens operations.

\begin{definition}[DoNothing Reduction]
    \label{def:donothing_reduction}
    The DoNothing reduction is a zero-round protocol with the following components:
    \begin{itemize}
        \item \textbf{Protocol specification}: $\pSpec := []$ (empty protocol, no messages exchanged)
        \item \textbf{Prover}: Simply stores the input statement and witness, and outputs them unchanged
        \item \textbf{Verifier}: Takes the input statement and outputs it directly
        \item \textbf{Input relation}: Any relation $R_{\mathsf{in}} : \StmtIn \to \WitIn \to \Prop$
        \item \textbf{Output relation}: The same relation $R_{\mathsf{out}} := R_{\mathsf{in}}$
    \end{itemize}
    \lean{DoNothing.reduction, DoNothing.oracleReduction}
\end{definition}

\begin{theorem}[DoNothing Perfect Completeness]
    The DoNothing reduction satisfies perfect completeness for any input relation.
    \lean{DoNothing.reduction_perfectCompleteness, DoNothing.oracleReduction_perfectCompleteness}
    \uses{def:donothing_reduction}
\end{theorem}

The oracle version of DoNothing handles oracle statements by passing them through unchanged as well. The prover receives both non-oracle and oracle statements as input, and outputs them in the same form to the verifier.

\subsection{Sending the Witness}

A fundamental building block in many proof systems is the ability for the prover to transmit witness information to the verifier. The SendWitness reduction provides this functionality in both direct and oracle settings.

\begin{definition}[SendWitness Reduction]
    \label{def:sendwitness_reduction}
    The SendWitness reduction is a one-round protocol where the prover sends the complete witness to the verifier:
    \begin{itemize}
        \item \textbf{Protocol specification}: $\pSpec := [(\PtoVdir, \WitIn)]$ (single message from prover to verifier)
        \item \textbf{Prover}: Sends the witness $\mathbbm{w}$ as its single message
        \item \textbf{Verifier}: Receives the witness and combines it with the input statement to form the output
        \item \textbf{Input relation}: $R_{\mathsf{in}} : \StmtIn \to \WitIn \to \Prop$
        \item \textbf{Output relation}: $R_{\mathsf{out}} : (\StmtIn \times \WitIn) \to \Unit \to \Prop$ defined by $((\mathsf{stmt}, \mathsf{wit}), ()) \mapsto R_{\mathsf{in}}(\mathsf{stmt}, \mathsf{wit})$
    \end{itemize}
    \lean{SendWitness.reduction}
    \uses{def:reduction}
\end{definition}

\begin{theorem}[SendWitness Perfect Completeness]
    The SendWitness reduction satisfies perfect completeness.
    \lean{SendWitness.reduction_completeness}
    \uses{def:sendwitness_reduction}
\end{theorem}

In the oracle setting, we consider two variants:

\begin{definition}[SendWitness Oracle Reduction]
    \label{def:sendwitness_oracle_reduction}
    The oracle version handles witnesses that are indexed families of types with oracle interfaces:
    \begin{itemize}
        \item \textbf{Witness type}: $\WitIn : \iota_w \to \Type$ where each $\WitIn(i)$ has an oracle interface
        \item \textbf{Protocol specification}: $\pSpec := [(\PtoVdir, \forall i, \WitIn(i))]$
        \item \textbf{Output oracle statements}: Combination of input oracle statements and the transmitted witness
    \end{itemize}
    % \lean{SendWitness.oracleReduction}
    \uses{def:oracle_reduction}
\end{definition}

\begin{definition}[SendSingleWitness Oracle Reduction]
    \label{def:sendsinglewitness_oracle_reduction}
    A specialized variant for a single witness with oracle interface:
    \begin{itemize}
        \item \textbf{Witness type}: $\WitIn : \Type$ with oracle interface
        \item \textbf{Protocol specification}: $\pSpec := [(\PtoVdir, \WitIn)]$
        \item \textbf{Conversion}: Implicitly converts to indexed family $\WitIn : \Fin(1) \to \Type$
    \end{itemize}
    \lean{SendSingleWitness.oracleReduction}
\end{definition}

\begin{theorem}[SendSingleWitness Perfect Completeness]
    The SendSingleWitness oracle reduction satisfies perfect completeness.
    \lean{SendSingleWitness.oracleReduction_completeness}
    \uses{def:sendsinglewitness_oracle_reduction}
\end{theorem}

\subsection{Random Query Testing}

One of the most fundamental oracle reductions tests whether two values computed from the statement and the oracle statements are equal, by comparing them at a single random query.

\begin{definition}[RandomQuery Oracle Reduction]
    \label{def:randomquery_oracle_reduction}
    Fix a type $M$ with an oracle interface and two maps $f, g$ that send a statement $x$ and oracle statements $\mathsf{o}$ to elements $f(x, \mathsf{o}), g(x, \mathsf{o}) \in M$. The RandomQuery reduction is:
    \begin{itemize}
        \item \textbf{Protocol specification}: $\pSpec := [(\VtoPdir, \mathsf{Query})]$ (single challenge from verifier)
        \item \textbf{Input relation}: $R_{\mathsf{in}}((x, \mathsf{o}), ()) := (f(x, \mathsf{o}) = g(x, \mathsf{o}))$
        \item \textbf{Verifier}: Samples a random query $q$ and sends it to the prover
        \item \textbf{Prover}: Receives $q$, performs no computation
        \item \textbf{Output relation}: $R_{\mathsf{out}}(((q, x), \mathsf{o}), ()) := (\mathsf{answer}(f(x, \mathsf{o}), q) = \mathsf{answer}(g(x, \mathsf{o}), q))$
    \end{itemize}
    Taking $f$ and $g$ to be two oracle statements gives the test of equality between two oracles; taking $g$ to be a value read from the statement checks a claimed query response.
    \lean{RandomQuery.oracleReduction}
    \uses{def:oracle_reduction}
\end{definition}

\begin{theorem}[RandomQuery Perfect Completeness]
    The RandomQuery oracle reduction satisfies perfect completeness: if $f(x, \mathsf{o}) = g(x, \mathsf{o})$, the two values agree on every query.
    \lean{RandomQuery.oracleReduction_completeness}
    \uses{def:randomquery_oracle_reduction}
\end{theorem}

\begin{theorem}[RandomQuery Knowledge Soundness]
    Suppose that whenever $f(x, \mathsf{o}) \neq g(x, \mathsf{o})$, a uniformly random query $q$ satisfies $\mathsf{answer}(f(x, \mathsf{o}), q) = \mathsf{answer}(g(x, \mathsf{o}), q)$ with probability at most $\varepsilon$. Then the RandomQuery oracle reduction satisfies worst-case round-by-round knowledge soundness with error $\varepsilon$.
    \lean{RandomQuery.oracleVerifier_rbrKnowledgeSoundnessWorstCase}
    \uses{def:randomquery_oracle_reduction}
\end{theorem}

The agreement bound only needs to hold for the pairs $f(x, \mathsf{o}), g(x, \mathsf{o})$ that actually occur. When it holds for all pairs, it is a property of the oracle interface:

\begin{definition}[Oracle Distance]
    \label{def:oracle_distance}
    For elements $a, b$ of $M$, their distance is
    \[ \mathsf{distance}(a, b) := |\{q : \mathsf{Query} \mid \mathsf{answer}(a, q) \neq \mathsf{answer}(b, q)\}|. \]
    The oracle interface has distance bound $d$ if every two distinct $a \neq b$ satisfy $\mathsf{distance}(a, b) \geq |\mathsf{Query}| - d$.
\end{definition}

\begin{theorem}[RandomQuery Knowledge Soundness from Distance]
    If the oracle interface on $M$ has distance bound $d$, then the RandomQuery oracle reduction satisfies round-by-round knowledge soundness with error $\frac{d}{|\mathsf{Query}|}$.
    \lean{RandomQuery.oracleVerifier_rbrKnowledgeSoundness}
    \uses{def:randomquery_oracle_reduction, def:oracle_distance}
\end{theorem}

\subsubsection{Random Evaluation of a Polynomial}

Polynomial identity testing is the instance of RandomQuery in which $M$ is the type of multivariate polynomials in $m$ variables over a field $\mathbb{F}$, a query is a point of $\mathbb{F}^m$ answered by evaluation, and $g = 0$. To test $P = Q$, take $f = P - Q$.

\begin{definition}[RandomEvaluation Oracle Reduction]
    \label{def:random_evaluation}
    Given a map $P$ from statements and oracle statements to polynomials in $\mathbb{F}[X_1, \ldots, X_m]$:
    \begin{itemize}
        \item \textbf{Input relation}: $P(x, \mathsf{o}) = 0$ as a polynomial
        \item \textbf{Protocol specification}: Single challenge $\tau \in \mathbb{F}^m$ from verifier to prover
        \item \textbf{Output relation}: $P(x, \mathsf{o})(\tau) = 0$ for the sampled point $\tau$
    \end{itemize}
    \lean{RandomEvaluation.oracleReduction}
    \uses{def:randomquery_oracle_reduction}
\end{definition}

\begin{theorem}[RandomEvaluation Soundness]
    If every $P(x, \mathsf{o})$ has total degree at most $d$, the RandomEvaluation reduction satisfies worst-case round-by-round knowledge soundness with error $\frac{d}{|\mathbb{F}|}$. This follows from the Schwartz--Zippel lemma: a nonzero polynomial of total degree at most $d$ vanishes at a uniformly random point of $\mathbb{F}^m$ with probability at most $\frac{d}{|\mathbb{F}|}$. For multilinear polynomials the error is $\frac{m}{|\mathbb{F}|}$.
    \lean{RandomEvaluation.oracleVerifier_rbrKnowledgeSoundnessWorstCase}
    \uses{def:random_evaluation}
\end{theorem}

\subsubsection{Batching Polynomial Evaluation Claims}

Another important building block is the ability to batch multiple polynomial evaluation claims into a single check using random linear combinations.

TODO: express this as a lifted version of $\mathsf{RandomQuery}$ over a virtual polynomial whose
variables are the random linear combination coefficients.

\begin{definition}[Batching Polynomial Evaluation Claims]
    \label{def:batching_polynomial_evaluation}
    Consider an $n$-tuple of values $v = (v_1, \ldots, v_n) \in \mathbb{F}^n$ and a polynomial map $E : \mathbb{F}^k \to \mathbb{F}^n$. The batching reduction is defined as:
    \begin{itemize}
        \item \textbf{Protocol specification}: Two messages:
        \begin{enumerate}
            \item Verifier sends random $r \in \mathbb{F}^k$ to prover
            \item Prover sends $\langle E(r), v \rangle := \sum_{i=1}^n E(r)_i \cdot v_i$ to verifier
        \end{enumerate}
        \item \textbf{Honest prover}: Computes the inner product $\langle E(r), v \rangle$ and sends it
        \item \textbf{Honest verifier}: Verifies that the received value equals the expected inner product
        \item \textbf{Extractor}: Trivial since there is no witness to extract
    \end{itemize}
\end{definition}

\begin{theorem}[Batching Completeness]
    The batching polynomial evaluation reduction satisfies perfect completeness.
\end{theorem}

\begin{remark}[Batching Security]
    The security of this reduction depends on the degree and non-degeneracy properties of the polynomial map $E$. The specific error bounds depend on the structure of $E$ and require careful analysis of the polynomial's properties.
\end{remark}

\subsection{Sending a Claim}

The SendClaim reduction enables a prover to transmit a claim (oracle statement) to the verifier, who then verifies a relationship between the original and transmitted claims.

\begin{definition}[SendClaim Oracle Reduction]
    \label{def:sendclaim_oracle_reduction}
    The SendClaim reduction is a one-round protocol for claim transmission:
    \begin{itemize}
        \item \textbf{Protocol specification}: $\pSpec := [(\PtoVdir, \OStmtIn)]$ (single oracle message)
        \item \textbf{Input}: Statement and single oracle statement (via $\mathsf{Unique}$ index type)
        \item \textbf{Prover}: Sends the input oracle statement as protocol message
        \item \textbf{Verifier}: Executes oracle computation $\mathsf{relComp} : \StmtIn \to \OracleComp[\OStmtIn]_{\mathcal{O}} \Unit$
        \item \textbf{Output oracle statements}: Sum type $\OStmtIn \oplus \OStmtIn$ containing both original and transmitted claims
        \item \textbf{Output relation}: $R_{\mathsf{out}}((), \mathsf{oracles}) := \mathsf{oracles}(\mathsf{inl}) = \mathsf{oracles}(\mathsf{inr})$
    \end{itemize}
    \lean{SendClaim.oracleReduction}
    \uses{def:oracle_reduction}
\end{definition}

\begin{theorem}[SendClaim Perfect Completeness]
    The SendClaim oracle reduction satisfies perfect completeness when the input relation matches the oracle computation requirement.
    % \lean{SendClaim.completeness} % TODO: theorem is not currently formalized.
    \uses{def:sendclaim_oracle_reduction}
\end{theorem}

\begin{remark}[SendClaim Development Status]
    The SendClaim reduction is currently under active development in the Lean formalization. Several components including the verifier embedding and completeness proof require further implementation. The current version represents a specialized case that may be generalized in future iterations.
    % TODO: Complete implementation as noted in the Lean code
\end{remark}

\subsection{Claim Reduction}

A fundamental building block for constructing complex proof systems is the ability to locally reduce one type of claim to another. The ReduceClaim reduction provides this functionality through mappings between statement and witness types.

\begin{definition}[ReduceClaim Reduction]
    \label{def:reduceclaim_reduction}
    The ReduceClaim reduction is a zero-round protocol that transforms claims via explicit mappings:
    \begin{itemize}
        \item \textbf{Protocol specification}: $\pSpec := []$ (no messages exchanged)
        \item \textbf{Statement mapping}: $\mathsf{mapStmt} : \StmtIn \to \StmtOut$
        \item \textbf{Witness mapping}: $\mathsf{mapWit} : \WitIn \to \WitOut$
        \item \textbf{Prover}: Applies both mappings to input statement and witness
        \item \textbf{Verifier}: Applies statement mapping to input statement
        \item \textbf{Input relation}: $R_{\mathsf{in}} : \StmtIn \to \WitIn \to \Prop$
        \item \textbf{Output relation}: $R_{\mathsf{out}} : \StmtOut \to \WitOut \to \Prop$
        \item \textbf{Relation condition}: $R_{\mathsf{in}}(\mathsf{stmt}, \mathsf{wit}) \iff R_{\mathsf{out}}(\mathsf{mapStmt}(\mathsf{stmt}), \mathsf{mapWit}(\mathsf{wit}))$
    \end{itemize}
    \lean{ReduceClaim.reduction}
    \uses{def:reduction}
\end{definition}

\begin{theorem}[ReduceClaim Perfect Completeness]
    The ReduceClaim reduction satisfies perfect completeness when the relation condition holds.
    \lean{ReduceClaim.reduction_completeness}
    \uses{def:reduceclaim_reduction}
\end{theorem}

\begin{definition}[ReduceClaim Oracle Reduction]
    \label{def:reduceclaim_oracle_reduction}
    The oracle version additionally handles oracle statements through an embedding:
    \begin{itemize}
        \item \textbf{Oracle statement mapping}: Embedding $\mathsf{embedIdx} : \iota_{\mathsf{out}} \hookrightarrow \iota_{\mathsf{in}}$
        \item \textbf{Type compatibility}: $\OStmtIn(\mathsf{embedIdx}(i)) = \OStmtOut(i)$ for all $i$
        \item \textbf{Oracle embedding}: Maps output oracle indices to corresponding input oracle indices
    \end{itemize}
    \lean{ReduceClaim.oracleReduction}
    \uses{def:oracle_reduction, def:reduceclaim_reduction}
\end{definition}

\begin{remark}[ReduceClaim Oracle Completeness]
    The oracle version's completeness proof is currently under development in the Lean formalization.
    % TODO: Complete when oracleReduction_completeness is implemented
\end{remark}

\subsection{Claim Verification}

Another essential building block is the ability to verify that a given predicate holds for a statement without requiring additional witness information.

\begin{definition}[CheckClaim Reduction]
    \label{def:checkclaim_reduction}
    The CheckClaim reduction is a zero-round protocol that verifies predicates:
    \begin{itemize}
        \item \textbf{Protocol specification}: $\pSpec := []$ (no messages exchanged)
        \item \textbf{Predicate}: $\mathsf{pred} : \StmtIn \to \Prop$ (decidable)
        \item \textbf{Prover}: Simply stores and outputs the input statement with unit witness
        \item \textbf{Verifier}: Checks $\mathsf{pred}(\mathsf{stmt})$ and outputs statement if successful
        \item \textbf{Input relation}: $R_{\mathsf{in}}(\mathsf{stmt}, ()) := \mathsf{pred}(\mathsf{stmt})$
        \item \textbf{Output relation}: $R_{\mathsf{out}}(\mathsf{stmt}, ()) := \mathsf{True}$ (trivial after verification)
    \end{itemize}
    \lean{CheckClaim.reduction}
    \uses{def:reduction}
\end{definition}

\begin{theorem}[CheckClaim Perfect Completeness]
    The CheckClaim reduction satisfies perfect completeness.
    \lean{CheckClaim.reduction_completeness}
    \uses{def:checkclaim_reduction}
\end{theorem}

\begin{definition}[CheckClaim Oracle Reduction]
    \label{def:checkclaim_oracle_reduction}
    The oracle version handles predicates that require oracle access:
    \begin{itemize}
        \item \textbf{Oracle predicate}: $\mathsf{pred} : \StmtIn \to \OracleComp[\OStmtIn]_{\mathcal{O}} \Prop$
        \item \textbf{Never-fails condition}: $\mathsf{pred}(\mathsf{stmt})$ never fails for any statement
        \item \textbf{Oracle computation}: Verifier executes oracle computation to check predicate
        \item \textbf{Input relation}: Defined via oracle simulation of the predicate
    \end{itemize}
    \lean{CheckClaim.oracleReduction}
    \uses{def:oracle_reduction, def:checkclaim_reduction}
\end{definition}

\begin{theorem}[CheckClaim Oracle Perfect Completeness]
    The CheckClaim oracle reduction satisfies perfect completeness.
    % \lean{CheckClaim.oracleReduction_completeness}  % TODO: theorem not yet stated in ArkLib (currently commented out in CheckClaim.lean); relink once proven
    \uses{def:checkclaim_oracle_reduction}
\end{theorem}

\begin{remark}[CheckClaim Security Analysis]
    The round-by-round knowledge soundness proofs for both reduction and oracle versions are currently under development in the Lean formalization.
\end{remark}
